\section*{Chapitre \ref{ch:s2:manipulative-strategies-and-how-they-are-caught} --- Stratégies de manipulation et moyens de les détecter}

\begin{solution}{exo:s2:manipulative-strategies-and-how-they-are-caught:1}
$24\,814 \times 0.005 \approx 124$ ordres exécutés~; les 24\,690 autres ont été annulés.
\end{solution}

\begin{solution}{exo:s2:manipulative-strategies-and-how-they-are-caught:2}
$0.01 \times 11\,500 = 115$ jours de compte légitimes signalés, contre 77 épisodes réels~: moins de la moitié des alertes sont vraies.
\end{solution}

\begin{solution}{exo:s2:manipulative-strategies-and-how-they-are-caught:3}
Environ 38.6~millions de dollars.
\end{solution}

\begin{solution}{exo:s2:manipulative-strategies-and-how-they-are-caught:4}
Annuler beaucoup est ce que font les teneurs de marché~: ils cotent en continu, remplacent leurs cotations à mesure que les cours bougent, et annulent plus que les spoofers
de l'échantillon. Ce qui distingue un spoofer, c'est d'annuler un type d'ordre tout en négociant un autre, ce que le ratio fait disparaître en moyenne.
\end{solution}

\begin{solution}{exo:s2:manipulative-strategies-and-how-they-are-caught:5}
Les fonds indiciels et les autres investisseurs doivent négocier à la clôture, car leurs indices de référence sont des cours de clôture~; un fort volume
à la clôture est leur comportement normal. Le dossier Optiver apporte l'essentiel~: une position valorisée au prix de règlement-livraison à l'opposé des transactions de clôture, et un mouvement qui s'inverse.
\end{solution}

\begin{solution}{exo:s2:manipulative-strategies-and-how-they-are-caught:6}
Presque tous les ordres sur les marchés modernes sont annulés, et les traders légitimes annulent quand les cours bougent ou après une exécution~; une règle fondée sur la seule annulation
interdirait la tenue de marché. Le texte de loi et le Market Abuse Regulation visent les ordres passés pour être annulés ou pour tromper, ce que les tribunaux déduisent
de la conception et des schémas, comme dans l'affaire Coscia.
\end{solution}

\begin{solution}{exo:s2:manipulative-strategies-and-how-they-are-caught:7}
63~\% à un taux de faux positifs de 0.1~\% (environ 11 alertes légitimes) et 100~\% à 10~\% (environ 1\,150). Capturer le dernier tiers des épisodes coûte cent
fois plus d'alertes~: le seuil est fixé par le nombre d'alertes que les enquêteurs peuvent lire.
\end{solution}

\begin{solution}{exo:s2:manipulative-strategies-and-how-they-are-caught:8}
Le ratio ordres sur transactions n'a détecté aucun des spoofers insérés et a signalé des teneurs de marché à la place. L'absence d'alerte d'un détecteur qui ne voit pas le
schéma ne dit rien de sa présence.
\end{solution}

\begin{solution}{pb:s2:manipulative-strategies-and-how-they-are-caught:1}
\begin{enumerate}
\item Le \omterm{def:s2:manipulative-strategies-and-how-they-are-caught:spoofing}{spoofing} consiste à afficher un cours acheteur ou vendeur dans l'intention de l'annuler avant exécution~; le \omterm{def:s2:manipulative-strategies-and-how-they-are-caught:layering}{layering} est un \omterm{def:s2:manipulative-strategies-and-how-they-are-caught:spoofing}{spoofing} portant sur plusieurs ordres à des niveaux successifs, maintenus
à l'écart du meilleur cours.
\item Petits et gros ordres de part et d'autre~; gros ordres annulés après un délai, une petite exécution ou une seule grosse exécution~; 24\,814 gros ordres dont
0.5~\% ont été exécutés, contre environ 52~\% des petits ordres~; plus de 450\,000 gros ordres et 1.4~millions de dollars.
\item Une plateforme modifiée pratiquant le \omterm{def:s2:manipulative-strategies-and-how-they-are-caught:layering}{layering}, avec quatre à six gros ordres de vente à trois ou quatre niveaux du meilleur cours vendeur, sur plus de 400 jours, pour plus de 40~millions de dollars de profits,
et près de 200~millions de dollars de pression avant le flash crash~; 38.6~millions de dollars par accord transactionnel en 2016.
\item \omterm{def:s2:manipulative-strategies-and-how-they-are-caught:spoofing}{Spoofing} au sens de 7~U.S.C.~6c(a)(5)(C)~; le règlement énumère les faux signaux, les négociations à la clôture, les ordres qui initient des tendances, la \omterm{def:s2:manipulative-strategies-and-how-they-are-caught:benchmark}{manipulation d'indice
de référence}, et des indicateurs tels que les ordres retirés avant exécution.
\item \omterm{def:s2:manipulative-strategies-and-how-they-are-caught:ignition}{Déclenchement de momentum} (momentum ignition)~: ordres destinés à lancer ou à exagérer une tendance~; \omterm{def:s2:manipulative-strategies-and-how-they-are-caught:close}{manipulation du cours de clôture} (marking the close)~: négociation importante à la clôture pour faire bouger un cours de clôture.
\item Des traders compensent de grandes positions Trading at Settlement en négociant en sens inverse pendant les dernières minutes de clôture~; 19 tentatives, 5 réussies~; 14~millions de dollars.
\item Fausses contributions ou négociations coordonnées pour faire bouger un indice de référence~; quatre banques ont plaidé coupables en 2015 de manipulation de fixings, avec des amendes de plus de 2.5~milliards de dollars.
\item Trois prétendus teneurs de marché et neuf personnes ont pratiqué du wash trading sur des crypto-actifs pour le compte de promoteurs, des robots générant des milliards de dollars de volume par
jour.
\item Teneurs de marché, rafraîchisseurs de cotations, fournisseurs de profondeur de carnet et traders directionnels, chacun déclenchant un détecteur~; spoofers insérés suivant le
schéma Coscia.
\item \omterm{def:s2:manipulative-strategies-and-how-they-are-caught:spoofing}{Spoofing}~: 0~\%, 37~\%, 51~\% et 77~\% détectés~; clôture~: 1~\%, 21~\%, 61~\%~; wash trading~: 74~\% et 97~\%, chacun à 1~\% de faux positifs.
\item Chaque caractéristique a des causes innocentes qui se produisent rarement ensemble~; le schéma du spoofer requiert les deux.
\item Point non tranché~: aucune décision de justice n'a été trouvée~; on l'évite en passant par des canaux de transaction privés.
\item À un taux de faux positifs de 1~\%, le détecteur combiné de \omterm{def:s2:manipulative-strategies-and-how-they-are-caught:spoofing}{spoofing} détecte 77~\% des cas et le ratio ordres sur transactions 0~\%, en signalant des teneurs de marché~; le détecteur de clôture
détecte 61~\% et signale 3.3~\% des fonds indiciels~; la part d'auto-appariements détecte 97~\% et signale 5.0~\% des sociétés multi-algorithmes.
\item L'alerte est classée et instruite~: les ordres qui la sous-tendent sont lus, et l'intention est recherchée dans la conception, les communications et les registres.
\item Sur des cas légitimes semblables, au seuil d'exploitation, en comptant les alertes par enquêteur, avec des épisodes issus de cas réels lorsque c'est possible.
\item Mélanger les tailles et les moments des ordres~; plus les gros ordres ressemblent à de vrais ordres, plus ils sont exécutés, ce qui est précisément le risque que la
manipulation était conçue pour éviter.
\item Ils décrivent les schémas observables qui définissent l'infraction et que les tribunaux ont retenus comme preuves de l'intention.
\item Le détecteur de surveillance des schémas d'annulation.
\item Les mark-outs valorisent l'information d'un client~; la surveillance cherche à savoir si les ordres d'un compte étaient destinés à être exécutés du tout.
\item Les ordres du teneur de marché sont destinés à être exécutés~; ceux du spoofer ne le sont pas.
\end{enumerate}
\end{solution}

\begin{solution}{iq:s2:manipulative-strategies-and-how-they-are-caught:1}
En montrant que les ordres sont authentiques~: ils s'exécutent à des taux comparables selon leur taille et leur côté, compte tenu de leur distance au meilleur cours, les annulations suivent les mouvements de cours
et les limites de risque plutôt que les exécutions du côté opposé, et les règles d'annulation du code ont des finalités légitimes, documentées avant toute négociation.
\end{solution}

\begin{solution}{iq:s2:manipulative-strategies-and-how-they-are-caught:2}
Par compte et par jour, comparer les taux d'exécution des gros et des petits ordres ainsi que le moment des annulations de gros ordres par rapport aux exécutions de l'autre côté~; agréger
leur combinaison~; tester sur des cas étiquetés et sur des stratégies légitimes choisies pour leur ressemblance~; rapporter le taux de détection au taux de faux
positifs d'exploitation.
\end{solution}

\begin{solution}{iq:s2:manipulative-strategies-and-how-they-are-caught:3}
Négocier pour des raisons réelles (couverture, ordres clients), sans chercher à faire bouger le règlement-livraison~; éviter de concentrer les négociations à contre-sens de la position dans
la fenêtre de clôture~; documenter les raisons~; vérifier les règles de la société et celles de la plateforme.
\end{solution}

\begin{solution}{iq:s2:manipulative-strategies-and-how-they-are-caught:4}
Revue préalable de la logique d'annulation et de taille des ordres des algorithmes~; prévention de l'auto-appariement~; surveillance des journaux d'ordres propres à la société avec les détecteurs
ci-dessus~; limites sur les négociations en fenêtre de clôture à contre-sens de positions valorisées au règlement-livraison~; surveillance des communications~; formation et escalade.
\end{solution}

\begin{solution}{iq:s2:manipulative-strategies-and-how-they-are-caught:5}
Chaque événement d'ordre (création, modification, annulation, exécution) avec horodatages de la bourse et locaux, compte et bénéficiaire effectif, algorithme et version, prix, taille,
sens et plateforme, ainsi que le motif des annulations lorsque le système le connaît~; conservés complets et non modifiés.
\end{solution}

\begin{solution}{iq:s2:manipulative-strategies-and-how-they-are-caught:6}
$77/(77 + 115) = 40.1~\%$ des alertes sont vraies. Avec dix épisodes au lieu de 100 et le même seuil, $7.7/(7.7 + 115) = 6.3~\%$~: des infractions plus rares rendent
les alertes de chaque détecteur majoritairement fausses.
\end{solution}
